\clearpage
\newcount\SavedTranslationPage
\SavedTranslationPage=\value{page}
\hypersetup{pageanchor=false}
\begingroup
\pagestyle{empty}\thispagestyle{empty}
\pdfbookmark[0]{Quellen und Terminologie der Übersetzungen}{translation-sources-I}
\noindent{\Large\bfseries Quellen und Terminologie der Übersetzungen\par}\medskip
\small\raggedright
Historischer Text: Carl Friedrich Gauss, \emph{Werke}, Band I, Göttingen, 1863. Historische Herausgeberbeiträge behalten ihre eigene Urheberschaft und ihre Unterschriften. Ursprünglich deutsche Passagen sind im deutschen Leseband in modernem Deutsch wiedergegeben; ihre historischen Originale werden gesondert mitgeliefert.

Diese englischen und deutschen Ausgaben sind mit OpenAI GPT-6.1 Sol KI-übersetzt; auch der Satz wurde durch KI erstellt. Die folgenden menschlichen Übersetzer und Autoren sind mit ihren Schriften Quellen für die jeweils bezeichnete Terminologie und den Sprachgebrauch. Mathematische Vorlage bleibt Gauss' originalsprachlicher Text.

\par\medskip\noindent\textbf{H. Maser.}\par\nobreak\smallskip\noindent Deutscher Übersetzer und Herausgeber, \emph{Carl Friedrich Gauss' Untersuchungen über höhere Arithmetik}, Berlin: Julius Springer, 1889. Entsprechende arithmetische Passagen und deutsche mathematische Terminologie.\par
\url{https://archive.org/details/carlfriedrichga00gausgoog}
\par\medskip\noindent\textbf{James Caddall Morehead und Adam Miller Hiltebeitel.}\par\nobreak\smallskip\noindent Übersetzer von Gauss' \emph{General Investigations of Curved Surfaces of 1827 and 1825}, Princeton University Library, 1902; Project Gutenberg 36856, TeX-Ausgabe. Verwandte englische geometrische Terminologie und Ausdrucksweise.\par
\url{https://www.gutenberg.org/ebooks/36856}
\par\medskip\noindent\textbf{Paul Taylor; Bernard Leak.}\par\nobreak\smallskip\noindent Taylors englische Übersetzung von Gauss' zweitem Beweis des Fundamentalsatzes der Algebra, mit Korrekturen von Leak, 1983. Verwandte englische algebraische Terminologie. Die herangezogene Online-Ausgabe nennt Dezember 1983 als Übersetzungsdatum, August 2003 für die LaTeX-Fassung und den 8. Juni 2007 als letztes Änderungsdatum der zugrunde liegenden Quelle. Ihre Danksagung nennt Mark Wainwright für die Rückgabe von Taylors Notizen und Nikita Danilov für die Korrektur von Tippfehlern.\par
\url{https://www.paultaylor.eu/misc/gauss.html}
\par\medskip\noindent\textbf{Ulrich Görtz.}\par\nobreak\smallskip\noindent \emph{Algebra}, WS 2021/22, 3.1.1–3.1.3 und Anfang von 3.1.4; \emph{Lineare Algebra II}, SS 2021, 18.4. Zeitgenössischer deutscher arithmetischer und algebraischer Sprachgebrauch.\par
\url{https://math.ug/a-ws2122/sect0006.html}\par
\url{https://math.ug/la2-ss21/sec-quotient-ring.html}
\par\medskip\noindent\textbf{Nikolay Moshchevitin.}\par\nobreak\smallskip\noindent \emph{Eine Bemerkung über positiv definite quadratische Formen und rationale Punkte}, arXiv:1510.06214v3. Verwandte zeitgenössische deutsche Terminologie quadratischer Formen.\par
\url{https://arxiv.org/abs/1510.06214v3}



\clearpage\endgroup\clearpage
\setcounter{page}{\the\SavedTranslationPage}
\hypersetup{pageanchor=true}
